\section{Product Scope}

SmartExam is built specifically for practical, computer-lab examinations \cite{scope1}. The system manages the entire process: pre-exam scheduling, student verification, live telemetry logging, secure file uploads, and post-exam grading reports. We organized the platform into twelve functional modules, starting with user login and hardware binding, and ending with audit logging.

We group the core features into five areas. First, we secure access by locking workstations to specific user accounts using unique Hardware IDs (HWIDs). Second, teachers set up exams with whitelisted applications and grading test cases. Third, during active exams, a background service monitors running processes and streams focus updates. Fourth, we auto-save student code and check for plagiarism. Lastly, we compile these metrics into downloadable PDF summaries.

Table~\ref{tab:Abbreviations} lists key abbreviations and definitions used in this document.

\begin{table}[!ht]
    \caption{Abbreviations}
    \label{tab:Abbreviations}
    \centering
    \begin{tabular}{ccl} 
        \hline
        \textbf{Serial \#} & \textbf{Abbreviation} & \textbf{Definition} \\
        \hline
        1 & AI & Artificial Intelligence \\
        2 & LMS & Learning Management System \\
        3 & CBT & Computer-Based Testing \\
        4 & SRS & Software Requirements Specification \\
        5 & RBAC & Role-Based Access Control \\
        6 & UI & User Interface \\
        7 & UX & User Experience \\
        8 & WBS & Work Breakdown Structure \\
        9 & CPM & Critical Path Method \\
        10 & QA & Quality Assurance \\
        11 & DBMS & Database Management System \\
        12 & API & Application Programming Interface \\
        13 & FR & Functional Requirement \\
        14 & NFR & Non-Functional Requirement \\
        \hline
    \end{tabular}
\end{table}

\subsection{Existing System Description}

Most universities still conduct practical lab exams manually. Teachers coordinate sessions over WhatsApp, mark attendance on paper, and collect final student code on USB drives. If a student is suspected of cheating, invigilators have to reconstruct the timeline from memory. This fragmented setup becomes chaotic as classes scale.

Below, we review ten popular testing tools, plagiarism checkers, and LMS platforms. While these systems solve individual problems, none offer a unified workflow. They lack the integration of desktop lockdown, live invigilation, automatic code execution, and evidence logs needed for physical lab exams.

\subsubsection{Respondus LockDown Browser}
Respondus LockDown Browser protects online test security by blocking other applications, secondary monitors, printing, and copy-paste functions during an exam session \cite{respondus1}. It is popular in universities because it plugs directly into existing learning management systems and can record webcam streams via Respondus Monitor. However, the system is designed for home testing and cannot allocate physical laboratory seats or bind students to specific lab computers. It also lacks a built-in compiler, meaning it cannot compile or grade code submissions, and it still requires external portals for file uploads.

\subsubsection{Safe Exam Browser}
Safe Exam Browser (SEB) is an open-source tool that locks down local workstations to run web-based assessments safely \cite{seb1}. It allows teachers to write custom configurations to whitelist specific web links and runs smoothly on Windows, macOS, and iOS. The drawback is that SEB does not handle student registration or seating maps. It also lacks coding compilation tools, has no built-in facial verification, and does not provide a live monitoring dashboard for invigilators, requiring external platforms for final submissions.

\subsubsection{MOSS}
Stanford's MOSS (Measure Of Software Similarity) is a widely used web service for catching plagiarism in coding assignments \cite{moss1}. It works by analyzing the structure of source code files across python, java, or C++ and automatically filters out starter template code. MOSS is highly efficient for batch checks, but it is not an exam management platform. It does not lock down workstations, monitor live class activity, verify student identities, or handle laboratory seat assignments.

\subsubsection{ExamSoft}
ExamSoft provides offline testing through a locked desktop client, making it a common choice for high-stakes professional assessments \cite{examsoft1}. Since it runs offline, students do not need a constant internet connection, and teachers get detailed feedback reports once data syncs. The main limitation is that it does not map physical lab rooms or track workstation details. It also lacks interactive compilers for coding exams, and because it runs offline, invigilators cannot monitor student screens or detect cheating in real-time.

\subsubsection{Moodle}
Moodle is an open-source learning management system (LMS) used to host courses, quizzes, and assignments online \cite{moodle1}. It manages user accounts and permissions well and supports a large library of community plugins. However, Moodle has no native way to lock down desktops or monitor background applications. It lacks a real-time supervision console for lab invigilators and does not store tamper-proof security logs for cheating disputes.

\subsubsection{Proctorio}
Proctorio is an automated proctoring platform that uses AI and webcams to flag suspicious behavior during online tests \cite{proctorio1}. It blocks browser tabs and integrates directly with existing university portals. However, the system is built for remote, home-based testing. It does not manage physical lab equipment, cannot assign seats based on local IP ranges, and relies on post-exam video flags rather than live invigator overrides.

\subsubsection{Honorlock}
Honorlock uses face-detection algorithms to verify student identities and prevent secondary device searches during online tests \cite{honorlock1}. While its browser locks and automated alerts work well for remote courses, the platform does not fit institutional lab exams. It lacks lab booking features, does not lock down the physical computer hardware, and cannot execute or compile code submissions.

\subsubsection{Blackboard Learn}
Blackboard Learn is a commercial LMS that provides standard quiz building, grading rubrics, and central student records \cite{blackboard1}. It has robust access control settings and connects with external educational plugins. However, the system is designed for general course administration. It lacks workstation lockdown tools, does not log background system processes, and requires third-party add-ons to secure lab-based exams.

\subsubsection{Canvas LMS}
Canvas LMS is a popular academic portal that handles course materials, online quizzes, and grading lists \cite{canvas1}. It supports easy plugin integration and allows teachers to manage roles. However, Canvas offers no workstation control or live telemetry. Teachers cannot monitor student screens, and the platform lacks scheduling tools for physical lab machines.

\subsubsection{Google Classroom}
Google Classroom is a lightweight, cloud-based platform for organizing assignments and sharing documents \cite{gclassroom1}. It connects easily with Google Drive and Meet, making it simple for teachers to distribute homework. However, it lacks secure testing features. It has no browser restrictions, cannot track workstation background apps, and offers no invigilation dashboard or machine mapping for lab rooms.

These reviews show that existing applications only cover specific parts of testing. Some lock the browser, while others check code similarity or manage online courses. A major gap remains. Currently, no single platform combines secure hardware binding, workstation lockdown, WebSockets telemetry, isolated code grading, and immutable audit logs in a single workflow. Table~\ref{tab:comparison} highlights these differences.

\begin{table}[H]
\centering
\caption{Applications Comparison}
\label{tab:comparison}
\renewcommand{\arraystretch}{1.3}
\small
\setlength{\tabcolsep}{3pt}
\begin{tabularx}{\textwidth}{|>{\raggedright\arraybackslash}X|C{0.6cm}|C{0.6cm}|C{0.6cm}|C{0.6cm}|C{0.6cm}|C{0.6cm}|C{0.6cm}|C{0.6cm}|C{0.6cm}|C{0.6cm}|C{1.0cm}|}
\hline\hline
\multicolumn{1}{|c|}{\multirow{2}{*}{\textbf{Features}}} & \multicolumn{11}{c|}{\textbf{Applications}} \\ \cline{2-12}
\multicolumn{1}{|c|}{} &
\rotatebox{90}{Respondus LockDown~\hspace{0.2cm}} &
\rotatebox{90}{Safe Exam Browser~\hspace{0.2cm}} &
\rotatebox{90}{MOSS~\hspace{0.2cm}} &
\rotatebox{90}{ExamSoft~\hspace{0.2cm}} &
\rotatebox{90}{Moodle~\hspace{0.2cm}} &
\rotatebox{90}{Proctorio~\hspace{0.2cm}} &
\rotatebox{90}{Honorlock~\hspace{0.2cm}} &
\rotatebox{90}{Blackboard Learn~\hspace{0.2cm}} &
\rotatebox{90}{Canvas LMS~\hspace{0.2cm}} &
\rotatebox{90}{Google Classroom~\hspace{0.2cm}} &
\rotatebox{90}{\textbf{Proposed System}\hspace{0.2cm}} \\
\hline\hline
Browser Lockdown        & \cmark & \cmark & \xmark & \cmark & \xmark & \cmark & \cmark & \xmark & \xmark & \xmark & \cmark \\ \hline
ID Verification         & \xmark & \xmark & \xmark & \cmark & \xmark & \cmark & \cmark & \xmark & \xmark & \xmark & \cmark \\ \hline
AI Monitoring           & \xmark & \xmark & \xmark & \xmark & \xmark & \cmark & \cmark & \xmark & \xmark & \xmark & \cmark \\ \hline
Exam Scheduling         & \xmark & \xmark & \xmark & \cmark & \cmark & \xmark & \xmark & \cmark & \cmark & \xmark & \cmark \\ \hline
Result Management       & \xmark & \xmark & \xmark & \cmark & \cmark & \xmark & \xmark & \cmark & \cmark & \xmark & \cmark \\ \hline
Lab/Workstation Control & \xmark & \xmark & \xmark & \xmark & \xmark & \xmark & \xmark & \xmark & \xmark & \xmark & \cmark \\
\hline\hline
\end{tabularx}
\end{table}

\subsection{Future System Usage Analysis}
We designed SmartExam to manage every stage of lab testing. Before the exam, admins set up roles, and teachers assign students to specific workstation seating maps. During the session, students run the exam on a locked client, while invigilators track focus timelines on a WebSockets dashboard and override locks if hardware issues arise. Once the exam ends, the system compiles grade sheets and generates PDF reports.

This platform is especially useful for practical courses like programming, database design, and networking. Because students write and run code on local machines, SmartExam standardizes exam environments and replaces paper coordination with automated workflows.

Lastly, the system guarantees accountability by saving all activity logs, focus-loss warnings, and manual teacher overrides in a relational database. This provides a clear, immutable audit trail to resolve grading disputes and verify exam integrity.